\def\stat{zac-ion}



\def\tit{НЕКОТОРЫЕ МЕТОДИЧЕСКИЕ ВОПРОСЫ ОЦЕНКИ УРОВНЯ ТЕХНОЛОГИЧЕСКОЙ ГОТОВНОСТИ 

ПРОЕКТОВ ИНФОРМАЦИОННЫХ СИСТЕМ}



\def\titkol{Некоторые методические вопросы оценки УТГ %уровня технологической готовности 

проектов информационных систем}



\def\aut{А.\,А.~Зацаринный$^1$, Ю.\,С.~Ионенков$^2$}



\def\autkol{А.\,А.~Зацаринный, Ю.\,С.~Ионенков}



\titel{\tit}{\aut}{\autkol}{\titkol}



\index{Зацаринный А.\,А.}

\index{Ионенков Ю.\,С.}

\index{Zatsarinny A.\,A.}

\index{Ionenkov Yu.\,S.}



%{\renewcommand{\thefootnote}{\fnsymbol{footnote}}

%\footnotetext[1]

%{Статья публикуется по представлению программного комитета VI Международной 

%в~МСЦ РАН в~рамках государственного задания по теме FNEF-2022-0014.}} 



\renewcommand{\thefootnote}{\arabic{footnote}}

\footnotetext[1]{Федеральный исследовательский центр <<Информатика и~управление>> Российской академии наук, 

AZatsarinny@ipiran.ru}

\footnotetext[2]{Федеральный исследовательский центр <<Информатика и~управление>> Российской академии наук, 

UIonenkov@ipiran.ru}



 \label{st\stat}



 \thispagestyle{headings}

 

 \vspace*{-12pt} 

 



\Abst{Статья посвящена вопросам оценки уровня технологической готовности (УТГ) проектов 

информационных систем (ИС). Представлен анализ существующих подходов к~оценке 

уровня готовности проектов за рубежом и~в России. Показано, что в~настоящее время все 

большее развитие получает подход, использующий методы оценки тех\-но\-ло\-ги\-че\-ской\!/\!про\-из\-вод\-ст\-вен\-ной\!/\!ры\-ноч\-ной готовности 

(technology\!/\!manufacturing\!/\!commercialization readiness levels~--- TRL\!/\!MRL\!/\!CRL), 

основанные на 9-уров\-не\-вой шкале и~позволяющие оценить уровень готовности 

технологии, производства и~готовности продукта к~выходу на рынок. Рассмотрены 

преимущества данных методов. Отмечено, что в~России они в~основном используются 

в~корпоративных проектах и~практически не применяются при разработке проектов ИС 

в~интересах государственного управления. Сформулированы предложения по 

совершенствованию методов оценки проектов ИС, включая методический подход 

к~оценке УТГ проектов.}



\KW{проект; информационная система; уровень готовности проекта; технологическая 

готовность; производственная готовность }



\DOI{10.14357\!/\!08696527220301}  



\vskip 10pt plus 9pt minus 6pt



 \thispagestyle{headings}

 

 \vspace*{-16pt}

 

\section{Введение}



 \vspace*{-2pt}



  Современный этап развития цивилизации характеризуется формированием 

нового технологического уклада~--- информационного общества, базовыми 

составляющими которого являются инфокоммуникационные технологии, 

активное использование информационных ресурсов, изменения в~структуре 

производства и~предоставления услуг. Существование информационного 

общества невозможно без эффективно функционирующих ИС. При этом ИС 

относятся к~классу сложных больших систем с~длительным сроком 

эксплуатации и~высокой стоимостью разработки. 

  

  Создание ИС, соответствующих современному уровню, требует не только 

наличия высококвалифицированных специалистов, современной технической 

и~технологической базы, но и~повышения роли нормативно-методического 

обеспечения для своевременного принятия обоснованных решений в~процессе 

проектирования и~производства систем. 

  

  Вопросам совершенствования методического аппарата оценки 

эффективности и~качества ИС посвящен ряд публикаций авторов, 

в~частности~[1, 2]. Наряду с~этим существенное значение приобретает выработка 

новых подходов к~оценке технологической и~производственной готовности 

инновационных разработок.

  

  В отечественной практике для оценки готовности проектов, как правило, 

используется подход, предполагающий выполнение и~приемку определенных 

стадий и~этапов работы в~соответствии с~ГОСТ (эскизный проект, технический 

проект, рабочая документация и~т.\,п.), который основывается на наличии 

документов, фиксирующих получение определенных результатов работ.

  

  Вместе с~тем в~настоящее время за рубежом для оценки уровня готовности 

проектов применяется метод оценки уровней  

тех\-но\-ло\-ги\-че\-ской\!/\!про\-из\-вод\-ст\-вен\-ной\!/\!ры\-ноч\-ной 

готовности проекта (TRL\!/\!MRL\!/\!CRL), 

использующий \mbox{9-уров}\-не\-вую шкалу~[3--5]. В России также начинает 

применяться подобный метод, основанный на использовании УТГ~[6]. Эти методы представляют собой 

метрику оценки готовности проектов на разных стадиях и~сопоставления 

различных проектов между собой. Уровни определяются по установленным 

правилам с~учетом концепции проекта, технических и~технологических 

требований. 

  

  В данной статье представлены общие подходы к~оценке уровня готовности  

на\-уч\-но-тех\-ни\-че\-ских проектов ИС в~России и~за рубежом, а~также 

сформулированы предложения по совершенствованию методов оценки 

технологической готовности проектов ИС с~учетом шкалы TRL.



\begin{table}\small %tabl1

%\vspace*{-12pt}

\begin{center}

\Caption{Уровни готовности проекта TRL/MRL/CRL}

\vspace*{2ex}



\tabcolsep=1pt

\begin{tabular}{|c|p{40mm}|p{38mm}|p{35mm}|}

\hline

\tabcolsep=0pt\begin{tabular}{c}Уровень\\ готов-\\ ности\end{tabular}&

\multicolumn{1}{c|}{\tabcolsep=0pt\begin{tabular}{c}Технологическая\\ готовность\\  (TRL)\end{tabular}}&

\multicolumn{1}{c|}{\tabcolsep=0pt\begin{tabular}{c}Производственная\\ готовность\\ (MRL)\end{tabular}}&

\multicolumn{1}{c|}{\tabcolsep=0pt\begin{tabular}{c}Рыночная\\  готовность\\ (CRL)\end{tabular}}\\

\hline

1&Сформулирована фундаментальная концепция, обоснована полезность технологии&Сделаны выводы по 

основным требованиям к~производству&Проведена оценка востребованности про\-дукта\\

\hline

2&Сформулирована техническая концепция технологии, установлены области применения.

Подготовлено ТЗ &Определена концепция производства&Оценены целевые потребительские 

сегменты, определены ключевые преимущества\\

\hline

3&Разработан макетный образец для демонстрации основных характеристик технологии&Подтверждена 

концепция производства &Проведен конкурентный анализ\\

\hline

4&Разработан детальный макет, проведены испытания базовых функций&Достигнута возможность 

изготовления технических средств в~лабораторных условиях&Определены поставщики и~партнеры, 

сформирована ценовая политика \\

\hline

5&Работоспособность технологии продемонстрирована на детальном макете&Достигнута возможность 

изготовления прототипов компонентов систем в~реальных условиях&Разработана бизнес-мо\-дель\\

\hline

6&Разработан полнофункциональный образец, подтверждены рабочие характеристики в~реальных 

условиях&Достигнута возможность изготовления прототипов систем и~подсистем с~использованием готовых 

элементов производства &Получена точная спецификация продукта, уточнена бизнес-мо\-дель\\

\hline

7&Прототип технологии продемонстрирован в~составе системы в~реальных условиях 

эксплуатации&Достигнута возможность изготовления систем и~подсистем в~условиях, близких 

к~реальным&Проведен предварительный вывод продукта на рынок\\

\hline

8&Сборка реального устройства и~демонстрация в~составе системы в~реальных условиях&Достигнута 

готовность к~началу мелкосерийного производства&Проанализированы результаты предварительного вывода 

продукта на рынок, проработаны замечания клиентов\\

\hline

9&Фактическое применение технологии в~окончательном виде в~реальных условиях&Начато мелкосерийное 

производство, под\-го\-тов\-ле\-на база для полномасштабного производства&Осуществлен вывод продукта на рынок\\

\hline

\end{tabular}

\end{center}

\end{table}





\section{Анализ существующих подходов к~оценке уровня готовности проектов информационных систем}



  В России в~настоящее время для оценки уровня готовности проектов 

преимущественно используется подход, основанный на выполнении 

определенных стадий и~этапов научно-исследовательских и~опыт\-но-кон\-ст\-рук\-тор\-ских работ (НИОКР) в~соответствии с~ГОСТ класса~34 

(ГОСТ 34.601-90) либо ГОСТ класса 15 (ГОСТ РВ 15.105-2001, ГОСТ РВ 

15.203-2001)~--- эскизный проект, технический проект, рабочая документация 

и~т.\,п. Этот подход основывается на документальной фиксации выполнения 

и~приемки определенных работ, предусмотренных ГОСТ (техническим заданием (ТЗ), планом 

выполнения работ и~т.\,п.) для каждой стадии (этапа) НИОКР.



  

  В то же время за рубежом и~в ряде организаций в~России все большее 

развитие получают методы оценки готовности проектов~--- 

 TRL\!/\!MRL\!/\!CRL~[3--5]. Эти методы используют 9-уров\-не\-вую шкалу, которая 

позволяет оценить уровень готовности технологии, производства и~готовности 

продукта к~выходу на рынок. Оценка выражается в~натуральных числах от~1 

до~9, при этом~9 соответствует наиболее высокому уровню. Шкала позволяет 

разработчикам и~заказчикам осуществлять контроль за процессом разработки: 

проект не перейдет на следующий уровень, пока не будет достигнут 

предыдущий.

  

  Уровни готовности проекта по основным составляющим: технологическая 

готовность, производственная готовность и~готовность к~выходу на рынок~--- 

представлены в~табл.~1~[7].

  

  Уровень готовности проекта присваивают на основе анализа его 

качественных и~количественных показателей. При этом должно обеспечиваться 

документальное подтверждение достигнутых показателей. Решение 

о~присвоении уровня принимают научно-исследовательские организации (НИО), 

занятые в~данной области. К~оценке также привлекаются независимые 

эксперты. 

  

  Уровень готовности проекта служит полезным инструментом, 

уста\-нав\-ли\-ва\-ющим согласованную терминологию и~единообразную 

методологию оценки уровня зрелости проектов. Данный метод представляет 

собой систематизированную метрику оценки продукта и~позволяет сравнивать 

уровни готовности различных продуктов. Он обладает гибкостью 

и~возможностью адаптации к~конкретному продукту, а также нуждам 

конкретных отраслей либо организаций. 

  

  





  В России получил распространение метод оценки уровня технологической 

готовности TRL, который апробирован крупнейшими технологическими 

лидерами (Росатом, Объединенная авиастроительная корпорация (ОАК), 

Объединенная двигателестроительная корпорация, фонд <<Сколково>> и~др.). 

Дирекцией на\-уч\-но-тех\-ни\-че\-ских программ Минобрнауки России на основе 

метода TRL разработаны две методологии оценки готовности  

на\-уч\-но-тех\-ни\-че\-ских проектов: методология комплексной оценки 

состояния на\-уч\-но-тех\-ни\-че\-ских проектов через уровень готовности 

технологий~--- TPRL (Technology Project Readiness Level)~[8] и~модель 

комплексной оценки технологической готовности инновационных  

на\-уч\-но-тех\-ни\-че\-ских проектов~[9]. Эти продукты рассчитаны на 

применение в~органах управления сложными комплексными проектами. Они 

позволяют выявлять динамику развития проектов, оценивать команды 

исследователей и~разработчиков, предсказывать неявные риски комплексных 

проектов, устранять возможные нарушения по срокам.

  

   Кроме того, Росимуществом разработаны <<Методические рекомендации по 

сопоставлению уровня технологического развития и~значений ключевых 

показателей эффективности акционерных обществ с~государственным 

участием, государственных корпораций, государственных компаний 

и~федеральных государственных унитарных предприятий с~уровнем развития 

и~показателями ведущих ком\-па\-ний-ана\-ло\-гов>>, предусматривающие 

использование шкал TRL и~MRL~[10]. 

  

  Близкий к~TRL метод, основанный на использовании 9-уров\-не\-вой шкалы, 

предусмотрен в~ГОСТ~Р 56861-2016~[6], который устанавливает УТГ

для проектов (табл.~2).



\begin{table}\small %tabl2

\begin{center}

\Caption{Уровни технологической готовности для проектов}

\vspace*{2ex}



\begin{tabular}{|c|p{97mm}|}

\hline

\tabcolsep=0pt\begin{tabular}{c}Уровень\\ технологической\\ готовности\end{tabular}&\multicolumn{1}{c|}{Примерный состав работ}\\

\hline

\multicolumn{1}{|c|}{\raisebox{-6pt}[0pt][0pt]{УТГ 1}}&Выявлены и~зафиксированы фундаментальные принципы технологии\\

\hline

УТГ 2&Концепция или выбор варианта\\

\hline

\multicolumn{1}{|c|}{\raisebox{-6pt}[0pt][0pt]{УТГ 3}}&Расчетное и~(или) экспериментальное обоснование эффективности технологий\\

\hline

\multicolumn{1}{|c|}{\raisebox{-6pt}[0pt][0pt]{УТГ 4}}&Исследование макетов и~(или) компонентов в~лабораторных условиях\\

\hline

\multicolumn{1}{|c|}{\raisebox{-6pt}[0pt][0pt]{УТГ 5}}&Верификация макетов и~(или) компонентов при подходящих условиях\\

\hline

\multicolumn{1}{|c|}{\raisebox{-6pt}[0pt][0pt]{УТГ 6}}&Моделирование систем (подсистем) или испытание трехмерных моделей в~подходящих 

условиях\\

\hline

\multicolumn{1}{|c|}{\raisebox{-6pt}[0pt][0pt]{УТГ 7}}&Разработка прототипа системы, продемонстрированная на действующем продукте\\

\hline

\multicolumn{1}{|c|}{\raisebox{-6pt}[0pt][0pt]{УТГ 8}}&Сборка реальной системы и~проверка работоспособности в~условиях, близких к~реальным\\

\hline

УТГ 9&Работа реальной системы в~реальных условиях\\

\hline

\end{tabular}

\end{center}

\vspace*{-12pt}

\end{table}



  Таким образом, широко распространенный за рубежом метод оценки 

го\-тов\-ности проекта TRL\!/\!MRL\!/\!CRL в~некоторой степени получает развитие 

и~в~России. При этом в~большей степени в~корпоративных проектах 

используются методы TRL и~MRL (Минобрнауки, Росатом, ОАК и~др.). Вместе с~тем эти методы практически не используются при разработке проектов ИС 

в~интересах государственного управления.

  

\section{Предложения по совершенствованию методов оценки проектов 

информационных систем}



  Долговременные проекты разработки инновационных технологий и~ИС с~их 

применением требуют существенных ресурсов и~связаны со значительными 

рисками. Начиная с~самых ранних стадий разработки таких систем, необходимо 

применять универсальные показатели, не зависящие от технологий, которые 

позволяют заказчикам и~разработчикам оценивать их состояние и~достигнутый 

прогресс в~сравнении с~запланированным.

  

  Метод, основанный на применении 9-уров\-не\-вой шкалы TRL\!/\!MRL\!/\!CRL, 

имеет ряд положительных сторон:

  \begin{itemize}

  \item  предполагает единый подход к~оценке статуса проектов;

  \item обеспечивает общее понимание статуса проекта;

  \item позволяет сравнивать уровни готовности различных продуктов;

  \item позволяет принимать решения о переходе проекта от стадии к~стадии;

  \item используется для принятия решений, касающихся финансирования 

проектов;

  \item позволяет выстраивать управление рисками;

  \item обеспечивает эффективное управление НИОКР (планирование, 

контроль).

  \end{itemize}

  

  Вместе с~тем, как отмечалось выше, этот метод практически не используется 

при разработке проектов ИС в~интересах государственного управления.

  

  Для адаптации данного метода к~использованию при разработке проектов ИС в~интересах государственного управления может быть предложено несколько 

групп мероприятий. 

  

  \textbf{Первая группа} направлена на создание нор\-ма\-тив\-но-пра\-во\-вой основы 

применения метода на федеральном уровне на базе постановлений 

и~распоряжений Правительства РФ, соответствующих ГОСТ, а на ведомственном 

уровне~--- ведомственных нормативных актов (приказы, распоряжения, 

методические указания и~т.\,п.).

  

  \textbf{Вторая группа} должна обеспечить увязку с~существующими стандартами 

разработки ИС (ГОСТ классов~15, 19, 34 и~др.), что позволит интегрировать 

подход, основанный на определении уровней готовности проектов, 

с~сущест\-ву\-ющи\-ми подходами к~разработке систем. Кроме того, это позволит 

формализовать требования к~составу и~структуре документов, подтверждающих 

достижение определенного уровня готовности проекта, и~разработать систему 

подтверждения уровней готовности. При этом следует учитывать, что, 

например, уровни готовности технологий TRL соответствуют различным 

стадиям (этапам), определенным в~ГОСТ, в~частности уровни готовности 

 TRL1--TRL3 соответствуют на\-уч\-но-ис\-сле\-до\-ва\-тель\-ским работам, TRL4--TRL6~--- стадии 

 опыт\-но-кон\-ст\-рук\-тор\-ских работ (ОКР),  

а~TRL7--TRL9~--- стадиям ОКР и~производства. В~соответствии с~этим 

целесообразно планирование этапов НИОКР по результатам с~этапами, 

кратными TRL (разработка технологии, разработка образца, готовность 

производства и~т.\,п.). 

  

 \textbf{Третья группа} охватывает комплекс работ по разработке критериев 

и~показателей оценки уровня готовности проектов. Эти показатели должны 

быть максимально формализованы, чтобы обеспечить однозначное понимание 

уровней готовности и~учитывать специфику конкретных типов ИС. При этом 

в~число показателей уровня готовности проектов обязательно должны быть 

включены показатели информационной безопасности. Применение в~ИС 

современных инфокоммуникационных технологий способствует и~появлению 

новых видов опасностей. В~этих условиях выполнение проектов 

в~существенной степени зависит от учета возможных уязвимостей и~соблюдения всех мер информационной безопас\-ности.

  

  \textbf{Четвертая группа} посвящена формированию системы оценки проектов, 

включающей как самостоятельную оценку, проводимую разработчиком, так 

и~оценку заказчиком (НИО заказчика) 

и~независимую экспертную оценку треть\-ей стороной. При этом самооценка, 

проводимая разработчиком, является первичной. Итоговая же оценка 

определяется рабочей группой, куда входят как представители разработчика, 

заказчика (НИО заказчика), так и~независимые эксперты. При этом 

целесообразно определить перечень организаций, из которых выбираются 

независимые эксперты.

  

  Мероприятия \textbf{пятой группы} направлены на совершенствование 

методического аппарата оценки УТГ проектов. 

Это совершенствование должно проводиться в~направлении разработки 

методик оценки уровня готовности на\-уч\-но-тех\-ни\-че\-ских проектов, 

содержащих перечень показателей уровня их готовности, а~также методов их 

расчета.

  

  Необходимо отметить, что в~методиках, изложенных в~работах~[8, 9], 

кроме показателей технологической готовности проекта, представленных 

в~табл.~1, предложен ряд дополнительных показателей, учитывающих риски, 

производственную, инженерную, организационную готовность и~т.\,п. Их 

разработка вызвана тем, что классическая шкала TRL не охватывает ряд 

аспектов, которые необходимо учитывать при оценке проекта в~целом, 

например производственную готовность, риски, сроки выполнения проекта 

и~т.\,п.

  

\section{Методический подход к~оценке уровня технологической 

готовности проектов}



  С учетом вышеизложенного может быть предложен следующий подход 

к~оценке УТГ проектов. Кроме показателей, 

характеризующих УТГ проекта (TRL1--TRL9), 

предлагается ввести ряд дополнительных обобщенных показателей, 

характеризующих уровень проекта по сравнению с~другими проектами, 

технологические, ресурсные, временн\'{ы}е и~другие его существующие 

и~потенциальные оценки: потенциальная эффективность результатов; 

реализуемость; стоимость и~сроки выполнения проекта; риски. Каждый из 

указанных показателей включает в~себя ряд частных показателей, которые 

оцениваются на основе вер\-баль\-но-чис\-ло\-вой шкалы (табл.~3). 





  Оценка УТГ по вер\-баль\-но-чис\-ло\-вой 

шкале проводится экспертами. После получения экспертных оценок проводится 

усреднение суждений экспертов по формуле

  $$

  a_{ij}= \sqrt[n]{a^1_{ij} a_{ij}^2 \cdots a_{ij}^n}\,,

  $$

где $a_{ij}^n$~--- оценки, выставленные соответствующими экспертами;  

$n$~--- число экспертов.



  Итоговая оценка УТГ проекта по 

дополнительным показателям может быть получена путем свертки оценок 

экспертов по группам~1--4 табл.~3 с~применением соответствующих весовых 

коэффициентов. При этом вначале рассчитываются обобщенные показатели по 

группам~1--3, а~затем итоговый показатель по всем четырем группам. Уровень 

технологической готовности проекта по дополнительным показателям 

позволяет оценить потенциальную эффективность его результатов, 

реализуемость и~риски, что облегчает принятие решений о~финансировании 

проекта и~его дальнейших перспективах. 



\section{Заключение}



  Реализация планов развития ИС предполагает определение показателей 

готовности соответствующих проектов. В~связи с~этим задача разработки 

методического аппарата, позволяющего проводить комплексную оценку 

готовности на\-уч\-но-тех\-ни\-че\-ских проектов, представляется весьма актуальной.



\begin{table}\small %tabl3

\begin{center}

\Caption{Дополнительные показатели для оценки УТГ проекта}

\vspace*{2ex}



\begin{tabular}{|p{40mm}|p{68mm}|}

\hline 

\multicolumn{1}{|c|}{\textbf{Показатели}}&\multicolumn{1}{c|}{\textbf{Вербально-числовая шкала}}\\

\hline

\multicolumn{2}{|c|}{1.\ Потенциальная эффективность результатов проекта}\\

\hline

Возможность улучшения характеристик технологии&1~--- низкая; 2~--- средняя; 3~--- 

высокая\\

\hline

Возможность решения новых задач&1~--- низкая; 2~--- средняя; 3~--- высокая\\

\hline

Уровень технологии по сравнению с~мировым&1~--- отставание; 2~--- соответствие; 3~--- 

опережение\\

\hline

\multicolumn{2}{|c|}{2. Реализуемость технологии}\\

\hline

Новизна технологии&1~--- основана на известных методах и~моделях; 

2~--- основана на усовершенствованных методах и~моделях; 3~--- основана на новых методах и~моделях\\

\hline

Обеспечение технологической базы&1~--- низкое; 2~--- среднее; 3~--- высокое\\

\hline

Обеспеченность кадрами&1~--- низкая; 2~--- средняя; 3~--- высокая\\

\hline

\multicolumn{2}{|c|}{3. Стоимость и~сроки выполнения проекта}\\

\hline

Оценка стоимости&1~--- значительно превышает заявленную; 2~--- незначительно превышает заявленную; 

3~--- соответствует заявленной\\

\hline

Оценка сроков&1~--- значительно превышают заявленные; 

2~--- незначительно превышают заявленные; 3~--- соответствуют заявленным\\

\hline

\multicolumn{2}{|c|}{4. Степень риска}\\

\hline

Оценка степени риска&1~--- значительный; 2~--- допустимый; 3~--- минимальный\\

\hline

\end{tabular}

\end{center}

\vspace*{-12pt}

\end{table}



  

  В настоящее время за рубежом и~в некоторой степени в~России для оценки 

уровня готовности проектов получил распространение метод, основанный на 

применении 9-уров\-не\-вой шкалы TRL\!/\!MRL\!/\!CRL и~позволяющий оценить 

уровень готовности технологии, производства и~готовности продукта к~выходу 

на рынок и~сравнить различные проекты между собой.

  

  В статье рассмотрены основные аспекты данного метода. Отмечено, что 

в~России он в~основном используется в~корпоративных проектах и~практически 

не применяется при разработке проектов ИС в~интересах государственного 

управ\-ле\-ния. Сформулированы предложения по совершенствованию методов 

оценки проектов ИС, включая методический подход к~оценке УТГ проектов, позволяющий получить количественные 

оценки. Эти оценки могут быть использованы для принятия различных 

управленческих решений, например для разработки пла\-на-гра\-фи\-ка работ, 

плана финансирования и~т.\,п.

  

{\small\frenchspacing

{%\baselineskip=10.8pt

\addcontentsline{toc}{section}{Литература}

\begin{thebibliography}{99}

  \bibitem{1-zo}

  \Au{Зацаринный А.\,А., Ионенков~Ю.\,С.} Оценка эффективности  

ин\-фор\-ма\-ци\-он\-но-те\-ле\-ком\-му\-ни\-ка\-ци\-он\-ных сис\-тем.~--- М.: НИПКЦ 

Восход-А, 2020. 120~с.

  \bibitem{2-zo}

  \Au{Зацаринный А.\,А., Ионенков~Ю.\,С.} Некоторые вопросы оценки качества 

информационных сис\-тем~/\!\!/ Системы и~средства информатики, 2021. Т.~31. №\,4.  

С.~4--17.

  \bibitem{3-zo}

   \Au{Mankins J.\,C.} Technology readiness level.~--- Advanced Concepts Office of Space 

Access and Technology NASA, 1995. {\sf 

https://aiaa.kavi.com/apps/group\_public/\linebreak download.php/2212/TRLs\_MankinsPaper\_1995.pdf}.

  \bibitem{4-zo}

  Manufacturing Readiness Level (MRL) Deskbook. Version 2.0.~--- OSD Manufacturing 

Technology Program, 2011. {\sf 

https://www.dodmrl.com/MRL\_Deskbook\_V2.pdf}.

  \bibitem{5-zo}

  Метрика <<Рыночная готовность CRL>>. {\sf 

https://www.rosatom.ru/upload/iblock/\linebreak fb2/fb23f5ee18bc97e7c0e2f14df9146f6a.pdf}.

  \bibitem{6-zo}

  ГОСТ Р 56861-2016. Система управления жизненным цик\-лом. Разработка концепции 

изделия и~технологий. Общие положения.~--- М.: Стандартинформ, 2017. 11~с.

  \bibitem{7-zo}

  Уровни готовности продукта/технологии по основным составляющим проекта: 

технологическая го\-тов\-ность, производственная го\-тов\-ность и~го\-тов\-ность к~выходу на рынок. 

{\sf http://ctt.volnc.ru/uploads/activity\_files/2020/08/14251.pdf}.

  \bibitem{8-zo}

   \Au{Петров А.\,Н., Сартори~А.\,В., Филимонов~А.\,В.} Комплексная оценка состояния 

на\-уч\-но-тех\-ни\-че\-ских проектов через уровень го\-тов\-ности технологий~/\!\!/ Экономика науки, 

2016. Т.~2. №\,4. С.~244--260.

  \bibitem{9-zo}

   \Au{Комаров А.\,В., Петров~А.\,Н., Сартори~А.\,В.} Модель комплексной оценки 

технологической готовности инновационных на\-уч\-но-тех\-ни\-че\-ских проектов~/\!\!/ 

Экономика науки, 2018. Т.~4. №\,1. С.~47--57.

  \bibitem{10-zo}

  Методические рекомендации по сопоставлению уровня технологического развития и~значений ключевых показателей эффективности акционерных обществ с~государственным 

участием, государственных корпораций, государственных компаний и~федеральных 

государственных унитарных предприятий с~уровнем развития и~показателями ведущих 

компаний-ана\-ло\-гов: Приложение №\,2 к~протоколу №\,2 от~19~сентября 2017~г.\ 

заседания Межведомственной рабочей группы по реализации приоритетов инновационного 

развития президиума Совета при Президенте РФ по модернизации и~инновационному 

развитию России. {\sf 

https://www.economy.\linebreak gov.ru/material/file/535256e8193016b438e70338a9a0d42d/metodic\_ta\_mrg.pdf}.



       \end{thebibliography}

} }







\vspace*{-3pt}



\hfill{\small\textit{Поступила в~редакцию 18.02.22}}



%\vspace*{8pt}



%\hrule



%\vspace*{2pt}



\newpage



\vspace*{-28pt}



%\hrule



%\vspace*{8pt}



\def\tit{SOME METHODOLOGICAL ISSUES\\ OF~ASSESSING THE~LEVEL OF~TECHNOLOGICAL 

READINESS\\ OF~INFORMATION SYSTEMS PROJECTS}







\def\titkol{Some methodological issues of~assessing the level of~technological 

readiness of~IS %information systems 

projects}



\def\aut{A.\,A.~Zatsarinny and Yu.\,S.~Ionenkov}



\def\autkol{A.\,A.~Zatsarinny and Yu.\,S.~Ionenkov}



\titel{\tit}{\aut}{\autkol}{\titkol}



\vspace*{-8pt}



\noindent

Federal Research Center ``Computer Science and Control'' of the Russian Academy of Sciences, 

44-2~Vavilov Str., Moscow 119333, Russian Federation



\def\leftfootline{\small{\textbf{\thepage}

\hfill Sistemy i Sredstva Informatiki~---

Systems and Means of Informatics 2022 vol~32 no\,3}

}%

 \def\rightfootline{\small{Sistemy i Sredstva Informatiki~---

 Systems and Means of Informatics 2022 vol~32 no\,3

\hfill \textbf{\thepage}}}



\vspace*{3pt}

  







\Abste{The article is devoted to the issues of assessing the level of technological readiness of 

information systems (IS) projects. General approaches to assessing the level of readiness of projects 

abroad and in Russia are presented. It is shown that at present, an approach using the methods of 

assessing technological/production/market readiness levels (TRL\!/\!MRL\!/\!CRL) based 

on a 9-level scale and allowing to assess the level of readiness of technology, production, and 

product readiness to enter the market is becoming increasingly developed. The advantages of these 

methods are considered. It is noted that in Russia, they are mainly used in corporate projects and are 

practically not used in the development of IS projects in the interests of public administration. 

Proposals are formulated to improve the methods of evaluating IS projects including 

a~methodological approach to assessing the level of technological readiness of projects.}



\KWE{project; information system; project readiness level; technological readiness; production 

readiness}



\DOI{10.14357\!/\!08696527220301} 



%\vspace*{-10pt}





% \Ack

 %\noindent

 





%\vspace*{-3pt}



\renewcommand{\bibname}{\protect\rmfamily References}

%\renewcommand{\bibname}{\large\protect\rm References}



{\small\frenchspacing

{ %\baselineskip=12pt

\addcontentsline{toc}{section}{References}

\begin{thebibliography}{99}

  \bibitem{1-zo-1}

\Aue{Zatsarinnyy, A.\,A., and Yu.\,S.~Ionenkov.} 2020. \textit{Otsenka effektivnosti 

informatsionno-telekommunikatsionnykh sistem} [Evaluation of the effectiveness of information 

and telecommunication systems]. Moscow: NIPKTs Voskhod-A. 120~p.

  \bibitem{2-zo-1}

\Aue{Zatsarinny, A.\,A., Yu.\,S.~Ionenkov, and A.\,P.~Suchkov.} 2021. Nekotorye voprosy otsenki 

kachestva informatsionnykh sistem [Some issues of information system quality assessment]. 

\textit{Sistemy i~Sredstva Informatiki~--- Systems and Means of Informatics} 31(4):4--17.

  \bibitem{3-zo-1}

\Aue{Mankins, J.\,C.} 1995. Technology readiness level. Advanced Concepts Office of Space 

Access and Technology NASA. Available at: {\sf 

https://aiaa.kavi.com/apps/group\_\linebreak public/download.php/2212/TRLs\_MankinsPaper\_1995.pdf} 

(accessed August~15, 2022).

  \bibitem{4-zo-1}

Manufacturing Readiness Level (MRL) Deskbook. Version 2.0. 2011.

OSD Manufacturing Technology Program. Available at: {\sf 

https://www.dodmrl.com/MRL\_\linebreak Deskbook\_V2.pdf}  (accessed August~15, 2022).

  \bibitem{5-zo-1}

Metrika ``Rynochnaya gotovnost' CRL'' [The ``Market readiness CRL'' metric]. Available at: {\sf 

https://www.rosatom.ru/upload/iblock/fb2/fb23f5ee18bc97e7c0e2\linebreak f14df9146f6a.pdf} (accessed 

August~15, 2022).

  \bibitem{6-zo-1}

GOST 56861-2016. 2017. Sistema upravleniya zhiznennym tsiklom. Razrabotka kontseptsii izdeliya 

i~tekhnologiy. Obshchie polozheniya [Lifecycle management system. Product concept and 

technology development. General provisions]. Moscow: Standartinform Publs. 11~p.

  \bibitem{7-zo-1}

Urovni gotovnosti produkta/tekhnologii po osnovnym sostavlyayushchim proekta: 

Tekhnologicheskaya gotovnost', proizvodstvennaya gotovnost' i~gotovnost' k~vykhodu na rynok 

[Product/technology readiness levels for the main components of the project: Technology readiness 

level/manufacturing readiness level/commercialization readiness level]. Available at: {\sf 

http://ctt.volnc.ru/uploads/activity\_files/2020/08/14251.\linebreak pdf} (accessed August~15, 2022).

  \bibitem{8-zo-1}

\Aue{Petrov, A.\,N., A.\,V.~Sartory, and A.\,V.~Filimonov.} 2016. Kompleksnaya otsenka 

sostoyaniya nauchno-tekhnicheskikh proektov cherez uroven' gotovnosti tekhnologiy 

[Comprehensive assessment of the status of scientific and technical projects using 

Technology Project Readiness Level]. \textit{Ekonomika nauki} [Economics of Science] 2(4):244--260.

  \bibitem{9-zo-1}

\Aue{Komarov, A.\,V., A.\,N.~Petrov, and A.\,V.~Sartory.} 2018. Model' kompleksnoy otsenki 

tekhnologicheskoy gotovnosti innovatsionnykh nauchno-tekhnicheskikh proektov [The model of 

integrated assessment of technological readiness of innovative scientific and technological projects]. 

\textit{Ekonomika nauki} [Economics of Science] 4(1):47--57.

  \bibitem{10-zo-1}

Metodicheskie rekomendatsii po sopostavleniyu urovnya tekhnologicheskogo razvitiya 

i~zna\-che\-niy klyu\-che\-vykh po\-ka\-za\-te\-ley ef\-fek\-tiv\-nosti ak\-tsio\-ner\-nykh obshchestv s~go\-su\-dar\-st\-ven\-nym 

uchastiem, gosudarstvennykh korporatsiy, go\-su\-dar\-st\-ven\-nykh kompaniy i~federal'nykh 

gosudarstvennykh unitarnykh predpriyatiy s~urov\-nem raz\-vi\-tiya i~po\-ka\-za\-te\-lya\-mi vedushchikh 

kompaniy-analogov:

Prilozhenie No.\,2 k~protokolu No.\,2 ot 19~sentyabrya 2017~g. 

zasedaniya Mezhvedomstvennoy rabochey gruppy po rea\-li\-za\-tsii prioritetov innovatsionnogo razvitiya prezidiuma Soveta 

pri Prezidente RF po mo\-der\-ni\-za\-tsii i~in\-no\-va\-tsi\-on\-no\-mu razvitiyu Rossii

[Methodological recommendations for comparing the level of technological 

development and the values of key performance indicators of state-owned joint-stock companies, 

state corporations, state-owned companies, and federal state unitary enterprises with the level of 

development and indicators of leading analog companies:

Appendix No.\,2 to Protocol No.\,2 dated September~19, 2017, of the Meeting of the Interdepartmental 

Working Group on the Implementation of Innovation Development Priorities of the Presidium of the Council under 

the President of the Russian Federation for Modernization and Innovative Development of Russia]. Available at: {\sf 

https://www.economy.gov.ru/material/file/535256\linebreak e8193016b438e70338a9a0d42d/metodic\_ta\_mrg.pdf} 

(accessed August~15, 2022).



\end{thebibliography}

} }



\vspace*{-3pt}



\hfill{\small\textit{Received February 18, 2022}} 



\vspace*{-14pt} 



\Contr



\vspace*{-3pt}



\noindent

\textbf{Zatsarinny Alexander A.} (b.\ 1951)~--- Doctor of Science in technology, professor, 

principal scientist, Federal Research Center ``Computer Science and Control'' of the Russian 

Academy of Sciences, 44-2~Vavilov Str., Moscow 119333, Russian Federation; 

\mbox{AZatsarinny@ipiran.ru}



\vspace*{3pt}



\noindent

\textbf{Ionenkov Yuriy S.} (b.\ 1956)~--- Candidate of Science (PhD) in technology, senior 

scientist, Federal Research Center ``Computer Science and Control'' of the Russian Academy of 

Sciences, 44-2~Vavilov Str., Moscow 119333, Russian Federation; \mbox{uionenkov@ipiran.ru}





  

 

\label{end\stat}



\renewcommand{\bibname}{\protect\rm Литература} 

